\section{Horizontal and Vertical Canonical Forms}

% Source-numbering note: arXiv:1606.00608 uses one theorem counter for
% definitions, propositions, theorems, examples, and observations in Section 4.
% The three examples after Theorem 4.9 are numbered 4.10--4.12; hence
% `Prop:IV.12` is Proposition 4.13 in the arXiv source.

\begin{definition}[Normalized BNT-refined horizontal form for an MPO]
    \label{def:horizontal_cf_mpdo}
    \lean{MPOTensor.IsHorizontalCF}
    \leanok
    \uses{def:mpo_tensor, def:sector_bnt_cf, def:gauge_equiv}
    An MPO tensor is in \emph{normalized BNT-refined horizontal form} if its
    doubled-index MPS tensor is a BNT sector decomposition with an independent
    gauge on every repeated copy:
    \begin{align}
        M^{ab}
         & =\bigoplus_{j=0}^{g-1}\bigoplus_{q=0}^{r_j-1}
        \mu_{j,q}X_{j,q}A_j^{ab}X_{j,q}^{-1}.
        \label{eq:mpdo_horizontal_cf}
    \end{align}
    Here the $A_j$ are the distinct minimal representatives, $r_j>0$, and
    every weight $\mu_{j,q}$ is nonzero, $|\mu_{j,q}|\leq1$, and at least
    one weight has modulus one. The BNT hypotheses include
    irreducibility, left-canonicality, normalized self-overlap, eventual
    linear independence of the representative MPV families, and the
    exclusion of similarity up to phase between distinct representatives.
    Repeated gauge-equivalent copies are retained over the same representative
    rather than separated as distinct blocks. The total bond dimension of
    the displayed direct sum is exactly $D$. Equivalently, with the
    block-diagonal matrix $X=\bigoplus_{j,q}X_{j,q}$,
    \begin{align}
        M^{ab}
         & =X\left[\bigoplus_{j,q}\mu_{j,q}A_j^{ab}\right]X^{-1},
        \notag                                                    \\
        X
         & =\bigoplus_{j,q}X_{j,q}.
        \notag
    \end{align}
    This is a strengthened representative-grouped form of the canonical
    decomposition in \cite[Section~2.3]{Cirac2016MPDO_arXiv}. In contrast with
    Definition~\ref{def:cpsv_canonical_form}, the retained direct sum has the
    full bond dimension, every copy weight is nonzero and normalized, and
    gauge-equivalent normal blocks are grouped over a BNT representative with
    a separate gauge for each copy. It is therefore not the literal CPSV
    canonical-form hypothesis.
    % Source anchors: CPSV16 lines 214--246 for literal CF and lines 281--301
    % for the BNT representative decomposition.
\end{definition}

\begin{definition}[MPV representation associated with horizontal canonical form]
    \label{def:horizontal_cf_mpv_representation}
    \lean{MPOTensor.HasHorizontalCFMPVRepresentation}
    \leanok
    \uses{def:mpo_tensor, def:sector_bnt_cf, def:same_mpv2_pos}
    An MPO tensor has a \emph{horizontal-canonical-form MPV representation} if
    its doubled-index MPS tensor generates the same positive-length MPV family as a
    BNT canonical form
    $\bigoplus_{j=0}^{g-1}\bigoplus_{q=0}^{r_j-1}\mu_{j,q}A_j$.
    Its length-$N$ matrix product vector is
    \begin{align}
        \sum_{j=0}^{g-1}
        \left(\sum_{q=0}^{r_j-1}\mu_{j,q}^{N}\right)V^{(N)}(A_j).
        \label{eq:mpdo_horizontal_mpv}
    \end{align}
    This is the MPV-level decomposition in
    \cite[Section~2.3, lines~298--301]{Cirac2016MPDO_arXiv}.
\end{definition}

\begin{theorem}[BNT-refined horizontal form gives the BNT MPV representation]
    \label{thm:horizontal_cf_gives_mpv_representation}
    \lean{MPOTensor.IsHorizontalCF.hasHorizontalCFMPVRepresentation}
    \leanok
    \uses{def:horizontal_cf_mpdo,
        def:horizontal_cf_mpv_representation}
    Every MPO tensor in normalized BNT-refined horizontal form has the
    associated positive-length BNT matrix-product-vector representation.
\end{theorem}

\begin{proof}
    \leanok
    Each copy similarity cancels inside every positive-length closed
    matrix-product trace.
\end{proof}

\begin{figure}[ht]
    \centering
    $K=X\left(\bigoplus_j\mathcal K_j\right)X^{-1}$
    \begin{center}
        % K = X ( bigoplus_j K_j ) X^{-1}  [CPSV16 eq. (CFK), lines 1660-1665]:
        % the horizontal canonical form of the doubled tensor K as a gauge
        % conjugation of the block direct sum of basis tensors.  Ink to
        % indices: each side exposes one virtual index at each end, an up
        % ket leg and a down bra leg (the paper draws the bra leg
        % dashed-up); the gauge capsules X and X^{-1} are matrices riding
        % the virtual wire and carry no physical leg, and the wires between
        % the capsules and the middle tensor are the contracted virtual
        % indices of the paper's chain.  The paper's block-index lines j
        % and q become the direct sum written in the middle label: the
        % weights mu_j are absorbed into the basis tensors, which depend on
        % the block index j only, while the gauge blocks X_{j,q} depend on
        % both indices.  Both sides carry the same boundary signature
        % (virtual-west=1 | virtual-east=1 | physical-up=1 | physical-down=1).
        \begin{tenkz}[physical=updown]
            \tn{K}
        \end{tenkz}
        $ = $
        \begin{tenkz}[physical=updown]
            \tnX{X} & \tn{\bigoplus_j\mathcal K_j} & \tnX{X^{-1}}
        \end{tenkz}
    \end{center}
    \caption{Horizontal canonical form of the doubled tensor: the gauge $X$,
        the direct sum of basis tensors $\mathcal K_j$ carrying the ket and bra
        indices, and $X^{-1}$. In the source
        specialized simple-biCF/ZCL display the weights $\mu_j$ are independent
        of the copy index and absorbed into the basis tensors, so the
        basis tensors depend only on the block index $j$ while the gauge factors
        $X_{j,q}$ depend on both block indices. This is
        \cite[Appendix~C.2, lines~1660--1665]{Cirac2016MPDO_arXiv}.}
    \label{fig:mpdo_horizontal_canonical_form}
\end{figure}

\begin{definition}[Vertically viewed tensor of an MPO]%
    \label{def:vertical_tensor_mpdo}
    \lean{MPOTensor.verticalTensor}
    \lean{MPOTensor.verticalTensor_finProdFinEquiv}
    \leanok
    \uses{def:mpo_tensor}
    The \emph{vertical direction} of an MPO tensor exchanges the notion of
    physical and virtual indices \cite[line~943]{Cirac2016MPDO_arXiv}: the
    vertically viewed tensor $\widetilde M$ has its letters indexed by the
    horizontal bond pair $(a,b)$, and each letter is the operator on the
    physical space with matrix elements
    $(\widetilde M_{ab})_{ij}=M^{ij}_{ab}$.
    Concatenating $\widetilde M$ vertically generates the matrix product
    vectors of the vertical direction.
\end{definition}

\begin{lemma}[Vertical reflection identity for the adjoint tensor]
    \label{lem:vertical_tensor_adjoint_reflection}
    \lean{MPOTensor.bondPairSwap}
    \lean{MPOTensor.bondPairSwap_finProdFinEquiv}
    \lean{MPOTensor.bondPairSwap_involutive}
    \lean{MPOTensor.verticalTensor_adjointTensor}
    \leanok
    \uses{def:vertical_tensor_mpdo, def:mpo_adjoint_tensor}
    Adjunction exchanges the two oriented horizontal bond indices in the
    vertical view:
    \begin{align}
        \widetilde{M^\dagger}_{a,b}
         & =(\widetilde M_{b,a})^\dagger.
        \label{eq:mpdo_adjoint_reflection}
    \end{align}
\end{lemma}

\begin{proof}
    \leanok
    This follows entrywise from the exchange of the ket and bra indices and
    conjugate transposition of the virtual matrix.
\end{proof}

\begin{definition}[Vertical canonical form for an MPO]\label{def:vertical_cf_mpdo}
    \lean{MPOTensor.IsVerticalCF}
    \leanok
    \uses{def:mpo_tensor, def:vertical_tensor_mpdo, def:bnt,
        def:block_diagonal_tensor}
    An MPO tensor is in \emph{vertical canonical form} if there are a basis of
    normal tensors $\{M_\alpha\}_\alpha$ for its vertically viewed tensor
    $\widetilde M$, positive diagonal matrices
    $\mu_\alpha=\diag(\omega_{\alpha,0},\ldots,\omega_{\alpha,r_\alpha-1})$
    with $r_\alpha\geq1$,
    and a matrix $U$ from the physical space onto the retained nonzero sector
    space such that
    \begin{align}
        U\widetilde MU^\dagger
         & =B,
        \notag                                           \\
        \widetilde M
         & =U^\dagger B U,
        \notag                                           \\
        B
         & =\bigoplus_\alpha \mu_\alpha\otimes M_\alpha.
        \label{eq:mpdo_vertical_cf}
    \end{align}
    This is the conclusion of
    \cite[Proposition~4.13, lines~1863--1870]{Cirac2016MPDO_arXiv}.
    In this orientation $UU^\dagger=I$ on the retained space, while
    $U^\dagger U$ is its support projection in the original physical space.
    The reconstruction identity says that every letter of $\widetilde M$ is
    supported on this projection, so the discarded orthogonal complement and
    the off-diagonal corners vanish.
    The orthogonal complement may be a zero sector: the source's general
    canonical-form construction explicitly permits the sum of the nonzero
    block dimensions to be strictly smaller than the original dimension
    \cite[lines~214--225]{Cirac2016MPDO_arXiv}.
    Since each $\mu_\alpha$ is diagonal, the summand
    $\mu_\alpha\otimes M_\alpha$ is the direct sum
    $\bigoplus_{k=0}^{r_\alpha-1}\omega_{\alpha,k}M_\alpha$ of weighted copies
    of $M_\alpha$, so the right-hand side is a block-diagonal tensor over the
    pairs $(\alpha,k)$ in which the multiplicities $r_\alpha$ and the diagonal
    entries $\omega_{\alpha,k}$ stay explicit; the coefficients
    $m_\alpha=\tr(\mu_\alpha)=\sum_{k=0}^{r_\alpha-1}\omega_{\alpha,k}$
    appearing in the renormalization fixed-point characterization are
    recovered from them.
\end{definition}

\begin{figure}[ht]
    \centering
    % U \widetilde M U^dagger = bigoplus_alpha mu_alpha (x) M_alpha
    % (arXiv:1606.00608, Proposition 4.13, eq. UMU and figure
    % MPDO_UMU.png).  The physical isometry does not represent a plain
    % Kronecker product: two copy-tensor junctions thread one shared
    % multiplicity index through both mu_alpha and M_alpha.  Only
    % M_alpha carries the open virtual legs.  Both sides have boundary
    % signature
    % (virtual-west=1 | virtual-east=1 | physical-up=3 |
    % physical-down=3).
    \begin{tenkz}[rows={op:none, op, op:none}, tensor style=box]
        \tn[pill, wide=3, legs at={1,2,3}]{U} & & \\
        \tn[wide=3]{\widetilde M} & & \\
        \tn[pill, wide=3, legs at={1,2,3}]{U^\dagger} & &
    \end{tenkz}
    $ = \bigoplus_\alpha $
    {\tenkzkernel
    \begin{tenkz}[rows={wire,wire}, cols=4, bonds=none]
        \tn[at=(1,1), name=d1,
            ports={90:physical, 270:physical, 0:physical}]{}
        \tn[at=(1,2), skin=box, name=mu,
            ports={180:physical, 0:physical}]{\mu_\alpha}
        \tn[at=(1,3), name=d2,
            ports={90:physical, 270:physical, 180:physical}]{}
        \tn[at={midway (1,3) and (2,3)}, name=d3,
            ports={90:physical, 270:physical, 0:physical}]{}
        \tn[at=(2,4), skin=box, name=M,
            ports={180:virtual, 0:virtual, 90:physical, 270:physical,
                   158:physical}]{M_\alpha}
        \tnwire{d1.90}{open n}
        \tnwire{d1.0}{mu.180}
        \tnwire{d1.270}{open s}
        \tnwire{d2.90}{open n}
        \tnwire{mu.0}{d2.180}
        \tnwire{d3.90}{d2.270}
        \tnwire{d3.270}{open s}
        \tnwire{d3.0}{M.158}
    \end{tenkz}}
    \caption{Vertical direct-sum form of the tensor: after the physical
        isometry $U$, each block factorizes as the positive diagonal weight
        $\mu_\alpha$ on the multiplicity factor times the basis tensor $M_\alpha$,
        and only $M_\alpha$ carries the virtual legs. This is the decomposition
        $U\widetilde MU^\dagger=\bigoplus_\alpha\mu_\alpha\otimes M_\alpha$ of
        \cite[Proposition~4.13, lines~1864--1870]{Cirac2016MPDO_arXiv}.}
    \label{fig:mpdo_vertical_direct_sum}
\end{figure}

\begin{remark}[Finite witnesses for block-injective canonical form]%
    \label{rem:word_tuple_span_top}
    \lean{MPSTensor.BlockEntryIndex}
    \lean{MPSTensor.wordEntryFamily}
    \lean{MPSTensor.duplicateScalarBlocks_counterexample}
    \lean{MPSTensor.wordTupleSpanTop_of_wordEntryFamily_linearIndependent}
    \lean{MPSTensor.hasBlockSelectorWords_of_wordTupleSpanTop}
    \lean{MPSTensor.hasBlockSelectorWords_of_wordEntryFamily_linearIndependent}
    \lean{MPSTensor.HasBiCF}
    \lean{MPSTensor.hasBiCF_of_wordTupleSpanTop}
    \lean{MPSTensor.hasBiCF_of_wordEntryFamily_linearIndependent}
    \lean{MPSTensor.HasBlockSelectorWords}
    \lean{MPSTensor.PropBlockInjective}
    \lean{MPSTensor.wordTupleSpanTop_of_propBlockInjective}
    \lean{MPSTensor.duplicateScalarBlocks_not_hasBiCF}
    \lean{MPOTensor.horizontalCFData_of_wordTupleSpanTop}
    \lean{MPOTensor.horizontalCFData_of_wordEntryFamily_linearIndependent}
    \lean{MPOTensor.horizontalCFData_of_propBlockInjective}
    \leanok
    Let $(A_k)_{k}$ be a finite family of block tensors of physical
    dimension~$d$, and fix a length $L$; write $A_k^w$ for the length-$L$
    word evaluation of $A_k$ along a word $w \in \{0,\ldots,d{-}1\}^L$. The
    finite-length spanning hypothesis is that the tuple-valued evaluations
    span the full product algebra of block matrices,
    \begin{align}
        \spn_{\C}\bigl\{\,(A_k^w)_k
        : w \in \{0,\ldots,d{-}1\}^L\,\bigr\}
         & =\textstyle\prod_k \MN{D_k}.
        \label{eq:mpdo_word_tuple_span}
    \end{align}
    The corresponding finite-length separation property is that some length
    $L$ makes the trace pairing against block words nondegenerate: there is
    a length $L$ such that, for every tuple $(\Delta_k)_k$ of block
    matrices,
    \begin{align}
        \Bigl(\forall\, w \in \{0,\ldots,d{-}1\}^L,\
        \textstyle\sum_k \tr(\Delta_k A_k^w) = 0\Bigr)
         & \implies \forall k,\ \Delta_k = 0.
        \label{eq:mpdo_finite_separation}
    \end{align}
    Spanning at length~$L$ gives separation at the same length: a trace
    functional that vanishes on every tuple $(A_k^w)_k$ vanishes on their
    span, hence on the whole product algebra, so nondegeneracy of the
    product trace pairing forces every $\Delta_k=0$. Combined with
    injectivity, left-canonicality, and nonzero block weights, this
    finite-length separation supplies the older per-copy trace-separation
    hypotheses for the block-diagonal tensor.

    The spanning hypothesis has an equivalent entrywise form. For a block
    index $k$ and matrix entry $(a,b)$, let
    $w\longmapsto(A_k^w)_{a,b}$
    be the corresponding scalar word family, indexed jointly by the triple
    $(k,a,b)$. Linear independence of this scalar family already implies
    the spanning condition, hence the separation property, and hence the
    older per-copy trace-separation hypotheses directly.

    The duplicate-scalar counterexample shows that the separation property
    is strictly stronger than blockwise injectivity, left-canonicality, and
    nonzero weights alone. Take $d=1$ and two blocks of bond dimension~$1$
    with $A_0^{0}=A_1^{0}=1$ and weights $\mu_0=1$, $\mu_1=2$: both blocks
    are injective and left-canonical, $(A_k^0)^\dagger A_k^0=1$, and the
    weights are nonzero, but at every length $L$ there is a unique word $w$
    and $A_0^w=A_1^w=1$, so $\Delta_0=1$, $\Delta_1=-1$ satisfies
    $\sum_k \tr(\Delta_k A_k^w)=0$ while $\Delta\neq0$. Thus the separation
    property fails, and no length makes the entrywise scalar family linearly
    independent.

    The block-injectivity proposition of
    \cite[lines~340--345]{Cirac2016MPDO_arXiv} gives a concrete sufficient
    criterion for the spanning condition: a common length $L$ with every
    block $L$-block-injective, together with a length $S$ at which, for
    every $k$, there is a coefficient family
    $(c_w)_{w\in\{0,\ldots,d{-}1\}^S}$ with
    \begin{align}
        \sum_w c_w A_k^w
         & =\Id,
        \notag   \\
        \sum_w c_w A_j^w
         & =0
        \quad(j\neq k).
        \notag
    \end{align}
    Concatenating the injective prefix with this finite selector family
    produces the spanning condition at length $L+S$, hence the separation
    property and the older per-copy trace-separation hypotheses.
\end{remark}

The inverse tensor of a block-injective canonical form satisfies
\begin{align}
    \sum_p(\mathcal K_j^p)_{\alpha,\beta}
    (K^{-1,p}_{j'})_{\alpha',\beta'}
     & =\delta_{\alpha,\alpha'}\delta_{\beta,\beta'}\delta_{j,j'}.
    \label{eq:mpdo_inverse_tensor}
\end{align}
Writing $K=X(\bigoplus_j\mathcal K_j)X^{-1}$ in this contraction places the
gauges on the virtual legs of $K$ and identifies the block index of
$\mathcal K_j$ with the label $j$:
\begin{tenkzequation}
    \tenkzkernel
    % K^{-1}K in block-injective canonical form (arXiv:1606.00608,
    % eq. biCFK, Appendix C.2 lines 1666--1670; MPDO_biCF.png is the
    % ink specification).  In panel 1 the vertical contraction joins the
    % physical input of K to the corresponding input of K^{-1}; the other
    % physical/output index of K^{-1} remains open above.  Its primitive
    % boundary signature is (virtual-west=2 | virtual-east=2 |
    % physical-up=1 | physical-down=0).  The paper distinguishes both
    % physical contractions by a dashed stroke; the kernel records the
    % contraction's type without a second stroke, so the join is the
    % ordinary bond.
    \begin{tenkz}[rows={op,op}, cols=1, west=open, east=open]
        \tn[skin=box, ports={90:physical}]{K^{-1}} \\
        \tn[skin=box]{K}
    \end{tenkz}
    $ = $
    % Panel 2 resolves each virtual boundary index into the tensor leg and
    % the j- and q-rails used by the block decomposition.  X and X^{-1}
    % tap both selector rails; the drawn tap descends to the j-dot and
    % continues through it to the q-dot, which reads identically on a
    % diagonal copy register.  K_j taps only the j-rail.  Thus the composite
    % virtual boundaries of panel 1 appear here as four primitive virtual
    % wires on each side (tensor, gauge, j, q); the open physical signature
    % remains (physical-up=1 | physical-down=0).
    \begin{tenkz}[rows={wire,wire,wire}, cols=3, bonds=none]
        \tn[at=(1,1), skin=box, name=X,
            ports={180:virtual, 0:virtual, 270:virtual}]{X}
        \tn[at=(1,2), skin=box, name=K,
            ports={180:virtual, 0:virtual, 90:physical,
                   270:virtual}]{\mathcal K_j}
        \tn[at=(1,3), skin=box, name=Xinv,
            ports={180:virtual, 0:virtual, 270:virtual}]{X^{-1}}
        \tn[at={1 n of (1,2)}, skin=box, name=Kinv,
            ports={180:virtual, 0:virtual, 90:physical,
                   270:physical}]{K^{-1}}
        \tn[at=(2,1), name=jX,
            ports={180:virtual, 0:virtual, 90:virtual, 270:virtual}]{}
        \tn[at=(2,2), name=jK, ports={180:virtual, 0:virtual, 90:virtual}]{}
        \tn[at=(2,3), name=jXinv,
            ports={180:virtual, 0:virtual, 90:virtual, 270:virtual}]{}
        \tn[at=(3,1), name=qX, ports={180:virtual, 0:virtual, 90:virtual}]{}
        \tn[at=(3,3), name=qXinv,
            ports={180:virtual, 0:virtual, 90:virtual}]{}
        \tnwire{X.180}{open w}
        \tnwire{X.0}{K.180}
        \tnwire{K.0}{Xinv.180}
        \tnwire{Xinv.0}{open e}
        \tnwire{Kinv.90}{open n}
        \tnwire{Kinv.270}{K.90}
        \tnwire{X.270}{jX.90}
        \tnwire{K.270}{jK.90}
        \tnwire{Xinv.270}{jXinv.90}
        \tnwire{jX.270}{qX.90}
        \tnwire{jXinv.270}{qXinv.90}
        \tnwire[name=jw]{jX.180}{open w}
        \tnwire{jX.0}{jK.180}
        \tnwire{jK.0}{jXinv.180}
        \tnwire{jXinv.0}{open e}
        \tnwire[name=qw]{qX.180}{open w}
        \tnwire{qX.0}{qXinv.180}
        \tnwire{qXinv.0}{open e}
        \tnmark[form=label, label pos=w]{on jw 0.8}{$j$}
        \tnmark[form=label, label pos=w]{on qw 0.8}{$q$}
    \end{tenkz}
    $ = $
    % Panel 3 is the componentwise identity.  The inward virtual legs of X
    % and X^{-1} are the two matrix-unit indices: each stands open as a
    % raised stub facing the centre, the left one on X's east side and the
    % right one on X^{-1}'s west side, where MPDO_biCF.png curls them into
    % hooks.
    % The vertical line is physical, not virtual: it is normal to the tensor
    % rail, continues the north/output port of K^{-1}, and ends at the
    % j-rail dot which records delta_{j,j'}.  The q-rail passes independently
    % below.  The primitive boundary signature is (virtual-west=4 |
    % virtual-east=4 | physical-up=1 | physical-down=0).
    \begin{tenkz}[rows={wire,wire,wire}, cols=3, bonds=none]
        \tn[at=(1,1), skin=box, name=X,
            ports={180:virtual, 45:virtual, 270:virtual}]{X}
        \tn[at=(1,3), skin=box, name=Xinv,
            ports={135:virtual, 0:virtual, 270:virtual}]{X^{-1}}
        \tn[at=(2,1), name=jX,
            ports={180:virtual, 0:virtual, 90:virtual, 270:virtual}]{}
        \tn[at=(2,2), name=delta,
            ports={180:virtual, 0:virtual, 90:physical}]{}
        \tn[at=(2,3), name=jXinv,
            ports={180:virtual, 0:virtual, 90:virtual, 270:virtual}]{}
        \tn[at=(3,1), name=qX, ports={180:virtual, 0:virtual, 90:virtual}]{}
        \tn[at=(3,3), name=qXinv,
            ports={180:virtual, 0:virtual, 90:virtual}]{}
        \tnwire{X.180}{open w}
        \tnwire{Xinv.0}{open e}
        \tnwire{delta.90}{open n}
        \tnwire{X.270}{jX.90}
        \tnwire{Xinv.270}{jXinv.90}
        \tnwire{jX.270}{qX.90}
        \tnwire{jXinv.270}{qXinv.90}
        \tnwire[name=jw]{jX.180}{open w}
        \tnwire{jX.0}{delta.180}
        \tnwire{delta.0}{jXinv.180}
        \tnwire{jXinv.0}{open e}
        \tnwire[name=qw]{qX.180}{open w}
        \tnwire{qX.0}{qXinv.180}
        \tnwire{qXinv.0}{open e}
        \tnmark[form=label, label pos=w]{on jw 0.8}{$j$}
        \tnmark[form=label, label pos=w]{on qw 0.8}{$q$}
    \end{tenkz}
\end{tenkzequation}
Consequently, $K^{-1}K=\bigoplus_j\Id_j$.
This is the inverse identity in
\cite[Appendix~C.2, lines~1666--1670]{Cirac2016MPDO_arXiv}.

\begin{theorem}[Finite block-separation criterion]
    \label{thm:propblockinj_abstract}
    \lean{MPSTensor.hasBiCF_of_propBlockInjective}
    \leanok
    \uses{rem:word_tuple_span_top}
    If a block family admits the finite-length data described in
    the block-injectivity proposition of
    \cite[lines~340--345]{Cirac2016MPDO_arXiv}, then it is in
    block-injective canonical form.
\end{theorem}

\begin{proof}
    \leanok
    \uses{rem:word_tuple_span_top}
    If the common injective prefix has length $L$ and the selector has length
    $S$, concatenation gives simultaneous words of length $L+S$. For every
    tuple $(X_j)_j$, the resulting coefficients satisfy
    $\sum_{|w|=L+S}c_w A_j^w=X_j$.
    Hence the simultaneous word tuples span the product algebra, and
    nondegeneracy of the trace pairing gives finite-length trace separation.
\end{proof}

\begin{lemma}[Single-block word-tuple span]
    \label{lem:single_block_word_tuple_span}
    \lean{MPSTensor.wordTupleSpanTop_of_card_eq_one_of_isNBlkInjective}
    \leanok
    \uses{def:l_blk_injective, rem:word_tuple_span_top}
    Let $(A_j)_{j=0}^{g-1}$ be a family with $g=1$. If its unique block is
    $L$-block injective, then the simultaneous length-$L$ word evaluations
    span the full one-block product algebra.
\end{lemma}

\begin{proof}
    \leanok
    \uses{def:l_blk_injective, rem:word_tuple_span_top}
    For every matrix $X$, block injectivity gives coefficients such that
    $X=\sum_{|w|=L}c_wA_0^w$.
    Since the family has only the block $A_0$, this is precisely the required
    simultaneous word-tuple spanning identity.
\end{proof}
